\documentclass{reportkit}

\setreportkitleftheader{REPORTKIT / v1.2 PRIMITIVE TEST}
\setreportkitfooter{New primitive acceptance test}
\setreportkitversion{v1.2.0}

\title{ReportKit v1.2 -- new diagram primitives}
\author{Test build}
\date{4 September 2026}

\begin{document}
\maketitle

\begin{execsummary}
This is an acceptance test for three new diagram primitives (\texttt{RKStack},
\texttt{RKCycle}, \texttt{RKFunnel}) and one new callout (\texttt{deliverablenote}),
added to support process/workflow guides -- content with accumulating evidence tiers,
genuinely cyclic loops, and staged conversion, none of which the existing matrix,
swimlane, network, or layer primitives express cleanly. Content is drawn from the
career-exploration framework this extension was built for.
\end{execsummary}

\section{Evidence stack}
Portfolio evidence is not all equally persuasive. A credential is common and weak;
a measured, externally validated outcome is rare and strong. \texttt{RKStack} makes
that asymmetry visible directly, rather than as a bulleted list where every line
reads with equal weight.

\begin{diagram}[
  type=stack,
  caption={The career evidence stack.},
  label={fig:stack},
  source={Illustrative progression; ordering reflects persuasiveness, not a measured population.},
  description={Six tiers rising from credential through measured impact, each narrower and rarer than the one below it.}
]
  \RKStackSetup{11}{1.05}{6}
  \RKStackTier{0}{Credential --- ``I studied X''}
  \RKStackTier{1}{Assessment --- scored well}
  \RKStackTier{2}{Project --- built something}
  \RKStackTier{3}{Demonstrated complexity --- deployed, monitored}
  \RKStackTier[rk accent]{4}{External validation --- others used it}
  \RKStackTier{5}{Measured impact --- quantified result}
\end{diagram}

\begin{deliverablenote}{Evidence Portfolio}
Students should be able to point to at least one artifact per tier they claim, not
just assert the tier verbally.
\end{deliverablenote}

\section{Continuous feedback loop}
The career-exploration process in Figure~\ref{fig:cycle} has no natural endpoint:
recruitment outcomes update the competency model, which changes what gets learned
next. A \texttt{sequence} chain of \texttt{RKNode}/\texttt{RKEdge} implies a
terminus; \texttt{RKCycle} does not.

\begin{diagram}[
  type=cycle,
  caption={The career exploration feedback loop.},
  label={fig:cycle},
  source={Illustrative conceptual cycle.},
  description={Six stages arranged in a ring, each connected to the next, looping back to the first.}
]
  \RKCycleSetup{4.3}
  \RKCycleNode{explore}{90}{Explore\\careers}
  \RKCycleNode{research}{30}{Research\\market}
  \RKCycleNode{assess}{-30}{Assess\\self}
  \RKCycleNode[rk accent]{build}{-90}{Build\\evidence}
  \RKCycleNode{recruit}{-150}{Enter\\recruitment}
  \RKCycleNode{update}{150}{Update\\hypothesis}

  \RKCycleEdge{explore}{research}
  \RKCycleEdge{research}{assess}
  \RKCycleEdge{assess}{build}
  \RKCycleEdge{build}{recruit}
  \RKCycleEdge{recruit}{update}
  \RKCycleEdge[repeat]{update}{explore}
\end{diagram}

\section{Recruitment funnel}
Where a student's conversion rate collapses is diagnostic: few interviews suggests a
targeting or resume problem, while technical passes followed by final-round rejection
suggests a fit or communication problem. \texttt{RKFunnel} makes the narrowing itself
the point.

\begin{redflag}{Illustrative, not measured}
This funnel's proportions are illustrative. If a report needs the shape driven by
real conversion counts, generate it with \texttt{reportkit\_viz.py}
(\texttt{rkv.bar\_chart} or a purpose-built funnel chart) instead, and keep this
native form for conceptual explanation only.
\end{redflag}

\begin{diagram}[
  type=funnel,
  caption={Recruitment conversion narrows at every stage.},
  label={fig:funnel},
  source={Illustrative shape; not derived from measured conversion data.},
  description={Four tapering stages: applications, screens, interviews, and offers.}
]
  \RKFunnelSetup{10}{1.15}
  \RKFunnelStage{0}{1.00}{0.74}{Applications}
  \RKFunnelStage{1}{0.74}{0.42}{Screens}
  \RKFunnelStage{2}{0.42}{0.20}{Interviews}
  \RKFunnelStage[rk accent]{3}{0.20}{0.10}{Offers}
\end{diagram}

\section{Grayscale and semantic check}
\begin{principle}{Semantics before decoration}
The new primitives follow the same rule as the existing ones: \texttt{RKStack} and
\texttt{RKFunnel} both narrow, but for opposite semantic reasons (rising rigor versus
falling population), so they must remain visually distinguishable without relying on
a shared visual trick that would blur that distinction.
\end{principle}

\end{document}
